\clearpage
\fancyhead[L]{\small JUNO.}
\hypertarget{native-265-lower}{}
\begin{center}
JUNO.

\rule{0.6\linewidth}{0.4pt}

Monatliche Correspondenz. Band XI. S. 184–193. 1805 Februar.

\rule{0.6\linewidth}{0.4pt}
\end{center}

—\quad—\quad—

Die interessantesten und wichtigsten Beiträge lieferte wie immer auch diesmal der mit ausdauerndem Fleisse fortarbeitende Dr. \textsc{Gauss}, und mit diesen fangen wir daher auch hier unsere Darstellung
der letzten Beobachtungen und Berechnungen, die Juno betreffend, an. Vergebens hatte dieser Astronom im Monat November 1804 auf Mittheilung guter Meridian-Beobachtungen der Juno gewartet, um
mittelst solcher fernere verbesserte Elemente der Juno liefern zu können; allein da er bis zum Decbr.
keine erhielt, und die aus den dritten Elementen berechneten Positionen einige Minuten von den beobachteten abwichen, so gründete er auf mehrere, theils selbst am Kreismikrometer gemachte, theils von
Dr. \textsc{Olbers} erhaltene Beobachtungen der Juno folgende IV. Elemente derselben.

\begin{center}
\emph{IV. Elemente der Juno-Bahn.}
\end{center}

\begin{center}
\begingroup\fontsize{9}{11}\selectfont\setlength{\tabcolsep}{3pt}\renewcommand{\arraystretch}{1.1}
\begin{tabular}{lr}
Epoche 1805 . . . . . . . . . & $42^\circ\ 41'\ 34''$ \\
Sonnenferne . . . . . . . . . & $233\ 23\ 47$ \\
aufsteigender Knoten . . . . . . & $171\ 4\ 12$ \\
Neigung . . . . . . . . . . . & $13\ 4\ 9$ \\
tägliche Bewegung . . . . . . . . & $812''091$ \\
Logarithmus der halben Achse . . . . & $0.426935$ \\
Excentricität . . . . . . . . . . & $0.256841$ \\
\end{tabular}\endgroup
\end{center}

\indent Nach diesen berechnete Dr. \textsc{Gauss} den geocentrischen Lauf derselben für die Monate Decbr. 1804,
Januar und Februar 1805. Die Ephemeride für den mittlern Monat theilten wir schon im vorigen Hefte
mit, und wir lassen daher hier nur die für den Decbr. und Febr. folgen [sind in diesem Werke nicht
abgedruckt.] —\quad—\quad—

Die von Dr. \textsc{Gauss} im Decbr. und Januar gemachten drei Beobachtungen waren folgende:

\begin{center}
\begingroup\fontsize{9}{11}\selectfont\setlength{\tabcolsep}{3pt}\renewcommand{\arraystretch}{1.1}
\begin{tabular}{l|r|r|r}
& mittl. Zeit & AR. $\JunoSymbol$ & \shortstack{Declinatio\\australis}\\ \hline
1804 December 29 & $6^{\mathrm u}\ 53^{\mathrm m}\ 43^{\mathrm s}$ & $7^\circ\ 23'\ 11''$ & $7^\circ\ 28'\ 2''$ \\
30 & $6\ 19\ 50$ & $7\ 44\ 5$ & $7\ 19\ 57$ \\
1805 Januar 5 & $5\ 6\ 57.5$ & $9\ 59\ 42$ & $6\ 30\ 56$ \\
\end{tabular}\endgroup
\end{center}
\par\Needspace{11\baselineskip}
\SourceLeaf{266}\par
\indent Er schrieb uns hierbei, dass ihm die letztere Beobachtung zwar an und für sich
gut scheine, allein dass sie durch die mit einem Stern aus der Histoire
céleste dabei angewandte Vergleichung sehr zweifelhaft werde. Dr. \textsc{Gauss}
verglich am 5. Januar die Juno achtmal mit dem Stern achter Grösse, dessen
Position in der Histoire céleste pag. 135,

\begin{center}
\begingroup\fontsize{10}{12}\selectfont\setlength{\tabcolsep}{3pt}\renewcommand{\arraystretch}{1.1}
\begin{tabular}{c|r|r}
& AR. & Zenith-Dist.\\ \hline
8 & $0^{\mathrm u}\ 40^{\mathrm m}\ 14^{\mathrm s}$ & $55^\circ\ 19'\ 15''$ \\
\end{tabular}\endgroup
\end{center}


\noindent angegeben ist. Allein schon nach dem blossen Augenmass schien es ihm, als
mache dieser Stern mit drei andern, rhomboidalisch darüber stehenden eine
andere Configuration, als aus den Angaben in der Histoire céleste folge. Da
Dr. \textsc{Gauss} eine genauere Ortsbestimmung dieses Sterns wünschte, so benutzten
wir hierzu den ersten heitern Abend, und fanden die Vermuthung des letztern
völlig gegründet, indem wir die Declin. nicht $6^\circ\,28'$, sondern
$6^\circ\,23'\,9''$ südlich erhielten, wonach denn die am 5. Januar
beobachtete Declination der Juno $6^\circ\,25'\,47''$ betragen würde. Aus den
IV. Elementen folgt für den

\begin{center}
\begingroup\fontsize{10}{12}\selectfont\setlength{\tabcolsep}{3pt}\renewcommand{\arraystretch}{1.1}
\begin{tabular}{l|r|r}
& AR. $\JunoSymbol$ & Declinat.\\ \hline
1804 December 30 & $7^\circ\ 44'\ 14''$ & $7^\circ\ 20'\ 16''$ \\
1805 Januar 5 & $9\ 59\ 44$ & $6\ 25\ 58$ \\
\end{tabular}\endgroup
\end{center}


\noindent womit die beobachteten Positionen nach der letztern Reduction sehr gut
harmoniren. Doch müssen wir bei unserer Bestimmung jenes zweifelhaften
Sterns bemerken, dass die grosse Helligkeit, die bei dessen Culmination noch
Statt fand, diese Beobachtungen unter die ganz zuverlässigen gerade nicht
zählen lässt. —\quad—\quad—

[Mit diesen IV. Elementen hat \textsc{Gauss} nach den oben Seite 106 abgedruckten
Vorschriften den Zodiacus der Juno berechnet, Mon. Corr. B. XI. S. 225–228.
1805 März.]

\hfill [\textsc{von Zach.}]

\clearpage
\fancyhead[L]{\small CERES.}
\begin{center}
CERES.
\end{center}

\begin{center}
Monatliche Correspondenz. Band XI. S. 283–292. 1805 März.
\end{center}

{\rightskip=0pt plus 1em
Auch bei der Ceres ist nun die Epoche ihrer Sichtbarkeit vorüber, und selbst
an grössern Aequatorial-Instrumenten dürfte ihre Beobachtung, da sie in
beträchtlicher Entfernung vom Meridian geschehen müsste, mit Schwierigkeiten
verknüpft sein. Wir müssen jetzt sechs bis sieben Monate darauf Verzicht
thun, diesen Wandelstern am Himmel zu verfolgen, und nur dahin streben, uns
bei der nächsten Wiedererscheinung seiner sogleich zu versichern. Dies wird
mittelst der von Dr. \textsc{Gauss} abermals verbesserten Elemente und danach
berechneten Ephemeride keine Schwierigkeit haben, da man jetzt die Bahn der
Ceres als sehr genau bestimmt annehmen kann, und wir liefern hier noch
alles, was künftig zu Erleichterung ihrer Beobachtung beitragen kann.\par}

Da die nach den IX. Elementen berechnete Ephemeride des geocentrischen Laufs
der Ceres im September und October ihre gerade Aufsteigung um neun Minuten
zu klein und die Abweichung um vier Minuten zu gross angab, so nahm Dr.
\textsc{Gauss} die drei im Jahr 1802, 1803 und 1804 erfolgten Oppositionen der Ceres
zu Hülfe und gründete darauf folgende X. Elemente
\SourceLeaf{267}
\begin{center}
\begingroup\fontsize{9}{11}\selectfont\setlength{\tabcolsep}{3pt}\renewcommand{\arraystretch}{1.1}
\begin{tabular}{lr}
Epoche Seeberger Meridian 1804 . . . . . . . . & $312^\circ\ 1'\ 33''5$ \\
Tägliche Bewegung . . . . . . . . . . . . . & $771''0524$ \\
Sonnenferne . . . . . . . . . . . . . . . & $326^\circ\ 26'\ 3''1$ \\
Excentricität . . . . . . . . . . . . . . . & $0.0784757$ \\
Logarithmus der halben Axe . . . . . . . . . & $0.04420004$ \\
$\AscendingNode$ 1804 . . . . . . . . . . . . . . . . . & $80^\circ\ 59'\ 12''$ \\
Neigung . . . . . . . . . . . . . . . . . & $10\ 37\ 45$ \\
\end{tabular}\endgroup
\end{center}

\indent Nach diesen berechnete Dr. \textsc{Gauss} den geocentrischen Lauf dieses Planeten vom 28. Juli 1805 bis
24. Mai 1806 [und gibt die gerade Aufsteigung so wie die Abweichung auf Minuten genau, den Abstand
von der Erde in drei Decimalstellen die Lichtstärke in fünf Decimalstellen für 12 Uhr Seeberger Zeit
von drei zu drei Tagen an.]

—\quad—\quad—

Er [\textsc{Gauss}] schrieb uns hierbei, dass \textsc{Bessel}, rühmlichst bekannt durch seine Abhandlung über
die Bahn des Cometen von 1607 (M. C. 1804 X. B. S. 425 f.) ihm bei Berechnung dieser Ephemeride dadurch
behülflich gewesen sei, dass er alle nöthige Sonnen-Örter dazu geliefert habe. Gewiss, jeder der
die Rechnungen kennt, die die Bestimmung der Elemente eines Planeten und dann jeder daraus herzuleitende
Ort erfordert, muss es bewundern, wie ein einzelner Mann in so kurzen Zeiträumen so vielfache
mühsame Rechnungen zu vollenden vermögend war. —\quad—\quad—

\hfill[\textsc{von Zach.}]
